\section{Predicting what the cameras will see}
\label{sec:prediction}

This section scores the predicted observations themselves, which is what the
two world models exist for. It first states the trap that makes a raw image
score misleading, then gives the measure that avoids it, and then shows the
predictions next to what happened. Table~\ref{tab:prediction} and
Figure~\ref{fig:predmargin} carry the numbers, and
Figure~\ref{fig:filmstrip} carries the pictures.

\paragraph{Repeating the present already scores well.}
A camera that barely changes is easy to predict: copying the last observed
frame forward matches the future closely. The fingertip cameras are the extreme
case. Their gels change little except at contact, so holding the last frame
already scores between about \num{21} and \num{27} decibels of peak
signal-to-noise ratio (PSNR) on them, against \num{14} to \num{19} on the
overhead camera.
A model reported without that baseline would look good on the fingertips for a
reason that has nothing to do with prediction.

\paragraph{The measure is the margin over holding.}
Every prediction is therefore scored beside the last observed frame held
still, and the verdict is the number of horizon steps on which the model beats
it by more than a tenth of a decibel. That is counted per camera and never
pooled, because four of the five cameras are gels where holding is strong and
the fifth is the one where the scene visibly moves; a pooled figure would be
dominated by whichever group made prediction easiest. The structural similarity
index (SSIM) is reported beside PSNR as a second opinion.

\begin{table}[t]
  \centering
  \caption{Future-observation prediction on the seven held-back recordings.
  \emph{Held} is the last observed frame repeated at every step. \emph{Steps
  won} counts the horizon steps, each \SI{0.8}{\second} apart, on which the
  model's PSNR exceeds holding by more than a tenth of a decibel. The two
  models' steps are not the same length: the small model's are
  \SI{0.8}{\second} apart, the large model's about \SI{0.13}{\second}, so
  step counts compare within a row and never between models. The large model
  reads the four fingertips as one composite image.}
  \label{tab:prediction}
  \small
  \begin{tabular}{llrrrrr}
\toprule
& & \multicolumn{2}{c}{PSNR (dB)} & \multicolumn{2}{c}{SSIM} & \\
\cmidrule(lr){3-4}\cmidrule(lr){5-6}
camera & model & model & held & model & held & steps won \\
\midrule
overhead & DreamZero & 15.5 & 14.1 & 0.494 & 0.432 & 5 of 6 \\
left arm left gripper & DreamZero & 30.1 & 25.9 & 0.895 & 0.857 & 5 of 6 \\
left arm right gripper & DreamZero & 28.0 & 22.9 & 0.861 & 0.759 & 5 of 6 \\
right arm left gripper & DreamZero & 24.8 & 20.7 & 0.788 & 0.688 & 5 of 6 \\
right arm right gripper & DreamZero & 26.8 & 23.6 & 0.869 & 0.812 & 6 of 6 \\
overhead & FastWAM & 18.5 & 18.6 & 0.711 & 0.695 & 2 of 8 \\
fingertips, composite & FastWAM & 28.9 & 27.2 & 0.940 & 0.913 & 6 of 8 \\
\bottomrule
\end{tabular}

\end{table}

\begin{figure}[t]
  \centering
  \includegraphics[width=\textwidth]{prediction_margin.png}
  \caption{Each world model's PSNR minus that of holding the last frame, per
  camera and per horizon step. Above the line the model predicts better than
  doing nothing, plotted against time rather than step because the two models
  predict over different spans: the small model from \SI{0.8}{\second} to
  \SI{4.8}{\second} ahead, the large one to about \SI{1}{\second}. For the
  small model the fingertip line is the mean of its four fingertip cameras; the
  large model is scored on its single composite.}
  \label{fig:predmargin}
\end{figure}

\textbf{The small world model predicts the future better than holding the
present, on every camera, from the second step onwards.} It wins five or six
of six steps on each camera, by roughly three to eight decibels beyond the
first step. The margin is largest on the fingertip cameras, which is where
holding is strongest, so the model is not winning where the baseline is weak.
It loses or ties only at the first step, \SI{0.8}{\second} ahead, where the
future still looks almost exactly like the present.

\begin{figure}[t]
  \centering
  \includegraphics[width=\textwidth]{dreamzero_filmstrip.png}
  \caption{One held-back moment, as the small world model predicted it. For each camera the rows are what
  happened, what the model predicted, and the last observed frame held still;
  columns are horizon steps \SI{0.8}{\second} apart. The images are the ones
  the numbers in Table~\ref{tab:prediction} were computed from.}
  \label{fig:filmstrip}
\end{figure}

Figure~\ref{fig:filmstrip} shows what those margins look like. In the last
observed frame the arms are down over the garment and one fingertip gel shows
the bright marks of contact. In what actually followed, the arms lift clear and
the contact marks vanish. \textbf{The model predicts the release rather than
the present:} its fingertip rows lose the contact marks as the real ones do,
and its overhead rows show the arms lifted. Its images are blurrier than the
real ones, which is where the remaining error lies, but the change it predicts
is the change that happened.

\textbf{The large model predicts only a second ahead, and within that second
it gains on holding only at the fingertips.} On the overhead camera it stays
within half a decibel of the held frame throughout, winning two of its eight
steps by the stated margin. On the fingertip composite it starts below holding
and climbs to about five decibels above it by the end of its horizon. Its
horizon is set by how it was trained, not by this evaluation: it predicts nine
frames spaced four apart, so it never looks further ahead than the small
model's first step.

\begin{figure}[t]
  \centering
  \includegraphics[width=\textwidth]{fastwam_filmstrip.png}
  \caption{The large world model's prediction across its horizon, from the
  same recording as Figure~\ref{fig:filmstrip}; columns are about
  \SI{0.13}{\second} apart. The fingertip row is the model's two-by-two
  composite of the four fingertip cameras. The last observed frame differs
  from the small model's by \SI{0.8}{\second}, because the two models place
  it differently relative to the horizon.}
  \label{fig:fastwamstrip}
\end{figure}

Figure~\ref{fig:fastwamstrip} shows what the overhead number hides. The
large model's frames are sharp where the small model's are blurred, and they
move correctly: the grippers enter the top of the overhead view at the pace
they really did. The margin over holding stays near zero regardless, because
the garment filling most of the frame does not move within a second, so
holding is already right about most pixels and the motion that distinguishes a
prediction occupies a sliver of the image. \textbf{An image score averaged over
the frame can rate a correct prediction of a small moving part as no better
than doing nothing.} The fingertip composite, where contact changes a large
share of the pixels, is where the model's margin shows.

The two models overlap at one instant, \SI{0.8}{\second} ahead, and there the
large model is ahead on both panels of Figure~\ref{fig:predmargin}. That is
the only direct comparison the data allow, and it is one point. It is
consistent with a pretrained video prior helping in the near term, and it says
nothing about the seconds beyond, where only the small model predicts at all
and where it keeps its margin. Whether the prior buys anything over a longer
horizon would need the large model trained to predict one.

Three limits bound this section. Both models sample, and every number here uses
one pinned seed; a different seed gives a different rollout, and the spread
across seeds was not measured for these checkpoints. The two models see the
fingertips differently, as four images or one composite, so their fingertip
margins are comparable in direction and not to the decibel. And predicting the
next few seconds of the cameras well is not the same as acting well:
Section~\ref{sec:wmactions} is where their actions are scored.
